\chapter{Structuring Comments}\label{sec:StructuringComments}
Structuring comments basically help with navigating the source code, and therefore, they are the most simple type of comments. 

First, we have the \bdq{headers} that we already described in chapter \tqfullvref{sec:ClassAndInterfaceDeclarations}\footnote{Lines are ending at column~80, but to avoid a line break here, the samples are shortened.}:
\begin{lstlisting}[numbers=left,stepnumber=3,firstnumber=2]
    /*------------------*\
====** Enum Declaration **========================… …===================
    \*------------------*/
    /*---------------*\
====** Inner Classes **===========================… …===================
    \*---------------*/
    /*-----------*\
====** Constants **===============================… …===================
    \*-----------*/
    /*------------*\
====** Attributes **==============================… …===================
    \*------------*/
    /*------------------------*\
====** Static Initialisations **==================… …===================
    \*------------------------*/
    /*--------------*\
====** Constructors **============================… …===================
    \*--------------*/
    /*---------*\
====** Methods **=================================… …===================
    \*---------*/
\end{lstlisting}

They are used to separate the parts of a class from each other. If a class or interface does not have a particular part, the assigned structuring comment must be omitted.

Of course you can decide to use another format for structuring comments in your projects, but you should not forgo it. These \bdq{headings} significantly facilitate navigation within the source code.

The other type of structuring comments are the \bdq{element end comments}, repeating the name of the element in a comment after the closing curly brace.

For a type, that comment is in a line of its own and contains the kind of the type:

\keeplisting{4}
\begin{lstlisting}
public final class MyClass
{
    …
}
//  class MyClass
\end{lstlisting}
\keeplisting{4}
\begin{lstlisting}
public final interface MyInterface
{
    …
}
//  interface MyInterface
\end{lstlisting}
\keeplisting{4}
\begin{lstlisting}
public final enum MyEnum
{
    …
}
//  enum MyEnum
\end{lstlisting}

For a method or a constructor, the comment follows on the same line as the closing curly brace:
\keeplisting{3}
\begin{lstlisting}
public MyClass()
{ 
    … 
}   //  MyClass()
\end{lstlisting}
\keeplisting{3}
\begin{lstlisting}
public final void MyMethod()
{ 
    … 
}   //  MyMethod()
\end{lstlisting}

%==============================================================================
% End of file!!
%==============================================================================
